\documentclass[conference]{IEEEtran}
\IEEEoverridecommandlockouts
\usepackage{cite}
\usepackage{amsmath,amssymb,amsfonts}
\usepackage{algorithmic}
\usepackage{graphicx}
\usepackage{textcomp}
\usepackage{xcolor}
\usepackage{booktabs}
\usepackage{multirow}
\usepackage{balance}
\usepackage{url}
\usepackage{microtype}

\def\BibTeX{{\rm B\kern-.05em{\sc i\kern-.025em b}\kern-.08em
    T\kern-.1667em\lower.7ex\hbox{E}\kern-.125emX}}

\begin{document}

\title{Multi-Objective Portfolio Optimization Using Deep Reinforcement Learning and Market Regime Conditioning\\
\thanks{This research was conducted as a B.Tech Major Project at the Department of Computer Science and Engineering, School of Computer Science and Engineering (SCSE), IILM University, Greater Noida, India (Academic Session 2026--2027).}
}

\author{
\IEEEauthorblockN{Yug Bhandari\IEEEauthorrefmark{1}, Nimisha Mishra\IEEEauthorrefmark{1}, Yojit Bhatt\IEEEauthorrefmark{1}, Jahanvi Jha\IEEEauthorrefmark{1}, Sunny Ranjan\IEEEauthorrefmark{1}, and Vishakha Agarwal\IEEEauthorrefmark{2}}
\IEEEauthorblockA{\IEEEauthorrefmark{1}Department of Computer Science and Engineering, School of Computer Science and Engineering\\
IILM University, Greater Noida, Uttar Pradesh, India\\
Enrollment / Roll Nos: 2410031103, 2410030123, 2410031104, 2410030454, 2410030457\\
Email: \{yug.bhandari.ug24, nimisha.mishra.ug24, yojit.bhatt.ug24, jahanvi.jha.ug24, sunny.ranjan.ug24\}@iilm.edu}
\IEEEauthorblockA{\IEEEauthorrefmark{2}Assistant Professor, Department of Computer Science and Engineering\\
IILM University, Greater Noida, Uttar Pradesh, India\\
Email: vishakha.agarwal@iilm.edu}
}

\maketitle

\begin{abstract}
Quantitative portfolio optimization in non-stationary financial markets remains an open challenge due to rapid shifts between structural regimes (bullish trends, market corrections, and high-volatility turbulence). Classical Mean-Variance Optimization (MVO) relies on static distributional assumptions that fail during systemic disruptions, while naive deep reinforcement learning (DRL) policies frequently suffer from destabilizing turnover and acute policy collapse when encountering out-of-distribution market dynamics. In this paper, we propose an end-to-end multi-objective portfolio optimization framework that conditions a continuous Proximal Policy Optimization (PPO) agent on deterministic market regimes. The environment is formulated as a Markov Decision Process (MDP) over a 19-dimensional continuous feature space encompassing rolling returns, volatility metrics, regime one-hot indicators, and current portfolio allocations, producing action vectors mapped to the probability simplex via a Softmax transformation. The optimization objective incorporates net portfolio return, a quadratic variance risk penalty, and an explicit turnover transaction cost penalty (10 basis points). We evaluate our model on a 10-year dataset (2016--2025, 2,514 trading days) comprising leading equity assets (AAPL, MSFT, GOOGL, AMZN, NVDA) with a rigorous out-of-sample test window spanning 2024--2025 (502 unseen trading days). The proposed Regime-Aware PPO achieves a cumulative return of 99.80\% (CAGR 41.64\%), a Sharpe ratio of 1.33, and a Calmar ratio of 1.28 net of transaction costs. In contrast, an unconditioned PPO ablation achieves only 54.54\% return (Sharpe 0.89), demonstrating that regime-state conditioning yields a 45.26\% absolute return enhancement and superior risk-adjusted preservation across volatile economic cycles.
\end{abstract}

\begin{IEEEkeywords}
Deep Reinforcement Learning, Portfolio Optimization, Proximal Policy Optimization (PPO), Market Regime Detection, Multi-Objective Optimization, Quantitative Finance, Financial Machine Learning.
\end{IEEEkeywords}

\section{Introduction}
Modern quantitative finance is founded upon the seminal Mean-Variance Optimization (MVO) framework introduced by Harry Markowitz in 1952 \cite{markowitz1952portfolio}. By framing asset allocation as a quadratic programming problem balancing expected portfolio return against covariance risk, MVO established the theoretical bedrock of modern portfolio management. However, in contemporary dynamic markets, MVO exhibits fundamental structural shortcomings. First, MVO assumes that asset returns follow a stationary Gaussian distribution, an assumption repeatedly violated by empirical financial time series characterized by fat tails, volatility clustering, and skewness \cite{mandelbrot1963variation, cont2001empirical}. Second, classical sample covariance estimators are notoriously susceptible to estimation noise, causing MVO solvers to act as ``error maximizers'' that concentrate capital into extreme, fragile allocations \cite{michaud1989markowitz}. Third, traditional optimization schemes cannot adaptively adjust rebalancing frequencies in response to market friction, incurring substantial transaction costs and turnover drag.

To address non-stationarity, Deep Reinforcement Learning (DRL) has emerged as a promising paradigm for sequential financial decision-making \cite{jiang2017deep, liang2018adversarial}. Unlike supervised learning models that predict next-period asset prices, DRL agents interact directly with simulated market environments, learning continuous control policies that map observed market states directly to optimal portfolio weights $\mathbf{w}_t \in \Delta^N$. By optimizing long-term cumulative discounted rewards, DRL agents can theoretically capture nonlinear multi-asset interdependencies, adapt to time-varying correlation structures, and internalize transaction frictions.

Despite these theoretical advantages, deploying DRL to algorithmic asset management presents three severe practical hurdles:
\begin{enumerate}
    \item \textbf{Regime-Blind Policy Degradation:} Financial markets alternate between distinct macroeconomic regimes---prolonged secular bull markets, sudden liquidity-driven bear markets, and sideways high-volatility environments. Unconditioned DRL policies trained across mixed regimes often overfit to dominant bullish trends and experience severe drawdowns during sudden volatility spikes.
    \item \textbf{Transaction Cost Friction and Churn:} Model-free continuous action agents frequently oscillate between portfolio weights on a daily basis, causing excessive turnover that severely erodes real-world performance after accounting for brokerage commissions, exchange fees, and bid-ask slippage.
    \item \textbf{Regulatory and Concentration Risks:} Theoretical unbounded simplex weights can allocate up to 100\% of capital to a single volatile asset, violating realistic institutional compliance thresholds (e.g., UCITS 5/10/40 concentration limits).
\end{enumerate}

To overcome these challenges, this paper presents a comprehensive, production-grade framework for \textbf{Multi-Objective Portfolio Optimization Using Deep Reinforcement Learning and Market Regime Conditioning}. The primary contributions of this work are summarized as follows:
\begin{itemize}
    \item \textbf{Regime-Conditioned MDP Formulation:} We design a 19-dimensional continuous observation space that combines trailing asset returns, rolling volatility, current portfolio drift, and a deterministic 3-state market regime indicator (Bull, Bear, High Volatility), allowing the policy to dynamically modulate its risk aversion based on macro conditions.
    \item \textbf{Multi-Objective Penalized Reward Function:} We construct a customized reward formulation that explicitly balances daily net return against both rolling portfolio covariance variance and quadratic rebalancing turnover penalties under a realistic 10 basis points (0.10\%) transaction cost schedule.
    \item \textbf{Rigorous Empirical Benchmarking \& Ablation:} Over a 10-year span (2016--2025, 2,514 trading days) and an unseen 502-day out-of-sample test period (2024--2025), we evaluate our model against Equal Weight ($1/N$), Buy \& Hold, Markowitz MVO, and an unconditioned PPO ablation. We prove that regime conditioning yields a +45.26\% return improvement and elevates the Sharpe ratio from 0.89 to 1.33.
    \item \textbf{Full-Stack Architecture with Java Execution Bridge:} We develop an integrated software ecosystem featuring a Gymnasium environment, a Stable-Baselines3 PPO agent, an enterprise-grade Java 17 virtual trading engine enforcing strict single-asset 35\% concentration limits, and a FastAPI/Streamlit analytics dashboard.
\end{itemize}

The remainder of this paper is structured as follows: Section II reviews related literature in portfolio optimization and financial DRL. Section III details the mathematical problem formulation, regime classification, and policy optimization methodology. Section IV outlines the experimental setup, data pipeline, and benchmark models. Section V provides empirical results, ablation analyses, and transaction cost discussions. Finally, Section VI concludes the paper with directions for future research.

\section{Related Work}

\subsection{Classical Portfolio Optimization and Factor Models}
Portfolio optimization began with Markowitz's Mean-Variance framework \cite{markowitz1952portfolio}, subsequently extended into the Capital Asset Pricing Model (CAPM) by Sharpe \cite{sharpe1964capital} and Lintner \cite{lintner1965valuation}. To resolve the sensitivity of MVO to sample estimation noise, Black and Litterman \cite{black1992global} integrated subjective investor views with market equilibrium returns using Bayesian updating. Later approaches introduced mathematical regularization, including Ledoit and Wolf's shrinkage covariance estimators \cite{ledoit2004well} and DeMiguel et al.'s empirical demonstration that simple Equal Weighting ($1/N$) frequently outperforms complex optimizing models out-of-sample due to zero estimation error \cite{demiguel2009optimal}.

\subsection{Deep Reinforcement Learning in Asset Allocation}
Early applications of reinforcement learning to financial trading utilized Recurrent Reinforcement Learning (RRL) to optimize Sharpe ratios directly \cite{moody2001learning}. With the advent of deep learning, Jiang et al. \cite{jiang2017deep} introduced the Ensemble Indeterministic Intelligence Ensemble (EIIE) framework utilizing Deep Deterministic Policy Gradients (DDPG) and Convolutional Neural Networks for cryptocurrency portfolio rebalancing. Liang et al. \cite{liang2018adversarial} compared DDPG, Proximal Policy Optimization (PPO), and Policy Gradient methods in stock markets, highlighting that policy-gradient methods provide enhanced stability compared to value-based Q-learning in continuous portfolio spaces. Ding et al. \cite{ding2020deep} investigated multi-agent architectures for equity portfolios, while Zhang et al. \cite{zhang2020deep} demonstrated that cost-aware reward functions are mandatory to prevent catastrophic capital depletion from churn.

\subsection{Market Regime Identification and Regime-Switching}
Financial markets exhibit non-linear structural transitions driven by business cycles, monetary policy, and sentiment shocks. Hamilton's Markov-Switching models \cite{hamilton1989new} formalized regime-dependent transitions, demonstrating that asset return volatility varies dramatically across states \cite{guidolin2011markov}. Ang and Bekaert \cite{ang2002regime} showed that dynamic asset allocation conditioned on bear-market volatility states significantly improves out-of-sample investor utility. In machine learning, researchers have explored Hidden Markov Models (HMMs) and clustering methods \cite{kritzman2012regime} to categorize latent market states. However, latent statistical clusters often lack intuitive interpretability and can introduce look-ahead bias if re-estimated naively. In this paper, we employ robust, causal trend-and-volatility heuristic rules that classify regimes instantaneously without forward information leakage.

\section{Problem Formulation and Methodology}

\subsection{Markov Decision Process (MDP) Formulation}
We model discrete-time portfolio allocation as an infinite-horizon discounted Markov Decision Process denoted by the tuple $(\mathcal{S}, \mathcal{A}, \mathcal{P}, \mathcal{R}, \gamma)$, where:
\begin{itemize}
    \item $\mathcal{S} \subset \mathbb{R}^{D}$ is the continuous state space of dimension $D = 19$.
    \item $\mathcal{A} \subset \mathbb{R}^{N}$ is the continuous action space representing raw asset preference logits across $N = 5$ equity assets.
    \item $\mathcal{P}: \mathcal{S} \times \mathcal{A} \times \mathcal{S} \rightarrow [0, 1]$ represents the unknown transition probability distribution governed by real-world financial market price movements.
    \item $\mathcal{R}: \mathcal{S} \times \mathcal{A} \rightarrow \mathbb{R}$ is the multi-objective reward function.
    \item $\gamma \in [0, 1)$ is the temporal discount factor, set to $\gamma = 0.99$.
\end{itemize}

At each trading day $t \in \{1, 2, \dots, T\}$, the agent observes the current state vector $\mathbf{s}_t$, outputs an action vector $\mathbf{a}_t$, transitions to state $\mathbf{s}_{t+1}$, and receives a scalar reward $R_t$.

\subsection{State Space Representation}
To provide the agent with a holistic view of asset momentum, market risk, and portfolio inertia, the state vector $\mathbf{s}_t \in \mathbb{R}^{19}$ is formulated as the concatenation of five distinct sub-vectors:
\begin{equation}
    \mathbf{s}_t = \big[ \mathbf{r}_t, \; \boldsymbol{\sigma}_t, \; \mathbf{z}_t, \; \mathbf{w}_t^{\text{drift}}, \; c_t \big]^\top
\end{equation}
where:
\begin{enumerate}
    \item $\mathbf{r}_t = [r_{1,t}, \dots, r_{N,t}]^\top \in \mathbb{R}^N$: The 1-day percentage returns of the $N$ constituent assets.
    \item $\boldsymbol{\sigma}_t = [\sigma_{1,t}, \dots, \sigma_{N,t}]^\top \in \mathbb{R}^N$: The 20-day annualized rolling volatility for each asset, capturing localized price instability.
    \item $\mathbf{z}_t \in \{0, 1\}^3$: A one-hot encoded vector representing the current macroeconomic regime:
    \begin{equation}
        \mathbf{z}_t = 
        \begin{cases}
            [1, 0, 0]^\top, & \text{if Regime}_t = \text{BULL} \\
            [0, 1, 0]^\top, & \text{if Regime}_t = \text{BEAR} \\
            [0, 0, 1]^\top, & \text{if Regime}_t = \text{HIGH\_VOLATILITY}
        \end{cases}
    \end{equation}
    \item $\mathbf{w}_t^{\text{drift}} \in \Delta^N$: The current drifted asset weights prior to rebalancing at time $t$.
    \item $c_t \in [0, 1]$: The current portfolio unallocated cash ratio ($c_t = 0$ in fully invested settings).
\end{enumerate}
For $N = 5$ assets, the dimension is $5 + 5 + 3 + 5 + 1 = 19$. In the ablation setting without regime conditioning, $\mathbf{z}_t$ is omitted, yielding $D = 16$.

\subsection{Continuous Action Space and Simplex Mapping}
In long-only portfolio management without leverage, the target capital allocation $\mathbf{w}_t = [w_{1,t}, \dots, w_{N,t}]^\top$ must reside on the standard $N$-dimensional probability simplex:
\begin{equation}
    \Delta^N = \left\{ \mathbf{w}_t \in \mathbb{R}^N \; \middle| \; \sum_{i=1}^N w_{i,t} = 1.0, \quad w_{i,t} \ge 0, \; \forall i \right\}
\end{equation}
The policy network outputs an unconstrained continuous action vector $\mathbf{a}_t \in [-5.0, 5.0]^N$. To guarantee that the generated portfolio weights strictly satisfy the budget and non-negativity constraints without numerical instability, we apply a numerically stabilized Softmax transformation:
\begin{equation}
    w_{i,t} = \frac{\exp(a_{i,t} - \max_k a_{k,t})}{\sum_{j=1}^N \exp(a_{j,t} - \max_k a_{k,t})}
\end{equation}

\subsection{Portfolio Drift and Turnover Dynamics}
Between rebalancing intervals $t-1$ and $t$, asset prices evolve according to daily price relative vector $\mathbf{y}_t = [y_{1,t}, \dots, y_{N,t}]^\top$, where $y_{i,t} = P_{i,t} / P_{i,t-1} = 1 + r_{i,t}$. Consequently, the portfolio allocation naturally drifts due to differential asset performance. The drifted weight of asset $i$ immediately before executing action $\mathbf{a}_t$ is:
\begin{equation}
    w_{i,t}^{\text{drift}} = \frac{w_{i,t-1} (1 + r_{i,t})}{\sum_{j=1}^N w_{j,t-1} (1 + r_{j,t})}
\end{equation}
The portfolio rebalancing turnover $\tau_t$ is defined as the total $L_1$-norm distance between the target rebalancing weights $\mathbf{w}_t$ and the pre-rebalancing drifted weights $\mathbf{w}_t^{\text{drift}}$:
\begin{equation}
    \tau_t = \sum_{i=1}^N \big| w_{i,t} - w_{i,t}^{\text{drift}} \big|
\end{equation}
Executing this reallocation incurs transaction fees and bid-ask slippage. Modeling an institutional transaction cost rate $c_{\text{fee}} = 0.0010$ (10 basis points), the total friction cost incurred is:
\begin{equation}
    C_t = c_{\text{fee}} \cdot \tau_t
\end{equation}

\subsection{Multi-Objective Penalized Reward Function}
Standard DRL formulations that maximize gross cumulative return $\sum_t r_t$ inevitably lead to hyperactive trading policies with unacceptable volatility and capital destruction from transaction fees. To enforce disciplined asset allocation, we formulate a multi-objective reward function balancing net financial return, localized covariance risk, and excessive portfolio turnover:
\begin{equation}
    R_t = 100 \cdot r_{\text{net}, t} - \lambda_{\text{risk}} \cdot 1000 \cdot \sigma_{\text{port}, t}^2 - \lambda_{\text{cost}} \cdot \tau_t
\end{equation}
where:
\begin{itemize}
    \item $r_{\text{net}, t} = \sum_{i=1}^N w_{i,t} r_{i,t} - C_t$ is the gross portfolio return net of transaction fees.
    \item $\sigma_{\text{port}, t}^2 = \mathbf{w}_t^\top \boldsymbol{\Sigma}_t \mathbf{w}_t$ is the estimated portfolio variance computed over a rolling lookback window of $K = 20$ trading days, where $\boldsymbol{\Sigma}_t = \text{Cov}(\mathbf{R}_{t-K:t})$.
    \item $\lambda_{\text{risk}} = 1.0$ is the risk-aversion hyperparameter penalizing allocation into volatile or co-dependent assets.
    \item $\lambda_{\text{cost}} = 0.5$ is the transaction cost dampener penalizing unneeded portfolio turnover.
    \item Scaling constants ($100$ and $1000$) normalize gradient magnitudes across return percentages and variance magnitudes for stable gradient descent.
\end{itemize}

\subsection{Market Regime Detection Framework}
To condition the policy on macro-structural trends without introducing future look-ahead bias, we design a three-state deterministic classification algorithm. Let $\bar{M}_{t, 20}$ denote the cross-sectional mean 20-day momentum of the asset universe, $\bar{S}_{t, 50}$ denote the cross-sectional mean price-to-50-day Simple Moving Average (SMA) ratio, and $\bar{\sigma}_{t, 20}$ denote the average 20-day annualized rolling volatility across all assets.

The regime state $Z_t \in \{\text{BULL}, \text{BEAR}, \text{HIGH\_VOLATILITY}\}$ is classified chronologically via the following priority rules:
\begin{enumerate}
    \item \textbf{High Volatility Priority State:}
    \begin{equation}
        \text{If } \bar{\sigma}_{t, 20} \ge Q_{0.75}(\bar{\sigma}_{20}) \implies Z_t = \text{HIGH\_VOLATILITY}
    \end{equation}
    where $Q_{0.75}(\bar{\sigma}_{20})$ is the 75th percentile historical volatility threshold ($0.3128$ in our dataset).
    \item \textbf{Bull Market State:}
    \begin{equation}
        \text{Else if } \bar{M}_{t, 20} > 0 \quad \text{and} \quad \bar{S}_{t, 50} > 1.0 \implies Z_t = \text{BULL}
    \end{equation}
    \item \textbf{Bear Market State:}
    \begin{equation}
        \text{Else if } \bar{M}_{t, 20} < 0 \quad \text{and} \quad \bar{S}_{t, 50} < 1.0 \implies Z_t = \text{BEAR}
    \end{equation}
    \item \textbf{Boundary Fallback:} Otherwise, $Z_t = \text{BULL}$ if $\bar{M}_{t, 20} \ge 0$, else $\text{BEAR}$.
\end{enumerate}

Across the 2,514-day dataset (2016--2025), this classifier partitions the market history into:
\begin{itemize}
    \item \textbf{BULL}: 1,484 trading days (59.0\%)
    \item \textbf{HIGH\_VOLATILITY}: 629 trading days (25.0\%)
    \item \textbf{BEAR}: 401 trading days (15.9\%)
\end{itemize}
This captures major historical drawdowns, including the 2018 tech selloff, the March 2020 COVID-19 crash, and the 2022 inflationary bear market.

\subsection{Proximal Policy Optimization (PPO)}
To train the allocation agent, we utilize Proximal Policy Optimization (PPO) \cite{schulman2017proximal}, an actor-critic algorithm that ensures stable policy updates by constraining policy drift via a clipped surrogate objective. 

Let $\pi_\theta(\mathbf{a}_t|\mathbf{s}_t)$ parameterize the policy with network parameters $\theta$, and let $r_t(\theta)$ denote the probability ratio between the current and previous policy:
\begin{equation}
    r_t(\theta) = \frac{\pi_\theta(\mathbf{a}_t | \mathbf{s}_t)}{\pi_{\theta_{\text{old}}}(\mathbf{a}_t | \mathbf{s}_t)}
\end{equation}
The clipped surrogate objective function is expressed as:
\begin{equation}
    L^{\text{CLIP}}(\theta) = \hat{\mathbb{E}}_t \left[ \min\left( r_t(\theta)\hat{A}_t, \; \text{clip}(r_t(\theta), 1-\epsilon, 1+\epsilon)\hat{A}_t \right) \right]
\end{equation}
where $\epsilon = 0.2$ is the clipping parameter and $\hat{A}_t$ is the Generalized Advantage Estimator (GAE) \cite{schulman2015high}:
\begin{equation}
    \hat{A}_t = \sum_{l=0}^{\infty} (\gamma \lambda)^l \delta_{t+l}^V
\end{equation}
with temporal difference residual $\delta_t^V = R_t + \gamma V_\phi(\mathbf{s}_{t+1}) - V_\phi(\mathbf{s}_t)$, discount factor $\gamma = 0.99$, and GAE factor $\lambda = 0.95$. The joint optimization objective incorporates the clipped policy loss, a mean-squared error value function loss $L^{VF}(\phi)$, and an entropy bonus $S[\pi_\theta]$ to encourage exploration:
\begin{equation}
    L(\theta, \phi) = -\hat{\mathbb{E}}_t \left[ L^{\text{CLIP}}(\theta) - c_1 \big( V_\phi(\mathbf{s}_t) - V_t^{\text{targ}} \big)^2 + c_2 S[\pi_\theta](\mathbf{s}_t) \right]
\end{equation}
where coefficients are set to $c_1 = 0.5$ and $c_2 = 0.01$.

\subsection{Post-Processing and Regulatory Risk Engine}
In institutional asset management, risk mandates restrict over-concentration in any individual stock. Following institutional standards (e.g., UCITS 5/10/40), we enforce a hard maximum concentration cap of $\omega_{\max} = 0.35$ (35\%) per equity asset. 

When the raw Softmax action produces weights $w_{i,t} > \omega_{\max}$, our deterministic risk engine caps the offending asset at $\omega_{\max}$ and iteratively redistributes the excess capital $\Delta w = w_{i,t} - \omega_{\max}$ across all unconstrained assets proportionally to their existing weights:
\begin{equation}
    w_{j,t}^{\text{post}} = w_{j,t} + \Delta w \cdot \frac{w_{j,t}}{\sum_{k \notin \mathcal{C}} w_{k,t}}, \quad \forall j \notin \mathcal{C}
\end{equation}
where $\mathcal{C} = \{i \mid w_{i,t} \ge \omega_{\max}\}$. This post-processing guarantees that $\sum_{i=1}^N w_{i,t}^{\text{post}} = 1.0$ and $w_{i,t}^{\text{post}} \le 0.35$ for all assets.

\section{Experimental Setup}

\subsection{Dataset and Preprocessing}
The empirical asset universe consists of five benchmark large-cap equities listed on NASDAQ: Apple Inc. (AAPL), Microsoft Corp. (MSFT), Alphabet Inc. (GOOGL), Amazon.com Inc. (AMZN), and NVIDIA Corp. (NVDA). Historical daily trading data was obtained from Yahoo Finance covering a 10-year period from January 4, 2016, through December 31, 2025, comprising 2,514 active trading days.

All asset prices were adjusted for stock splits, reverse splits, and cash dividend distributions to avoid artificial return discontinuities. Forward and backward imputation was applied to ensure uniform calendar alignment across all assets.

\subsection{Chronological Dataset Splitting}
To strictly eliminate look-ahead bias and simulate realistic out-of-sample forward trading, the dataset was chronologically partitioned into three non-overlapping subsets:
\begin{itemize}
    \item \textbf{Training Set:} January 4, 2016 to December 31, 2022 (1,563 trading days, $\sim$7 years). Used for feature calculation, regime definition, and policy gradient training.
    \item \textbf{Validation Set:} January 3, 2023 to December 29, 2023 (250 trading days, 1 year). Used for early stopping and hyperparameter calibration.
    \item \textbf{Out-of-Sample Test Set:} January 2, 2024 to December 31, 2025 (502 trading days, 2 full years). Kept completely isolated during model training and used exclusively for final benchmark evaluation.
\end{itemize}

\subsection{Benchmark Strategies}
We evaluate the proposed Regime-Aware PPO agent against four classical and reinforcement learning benchmarks:
\begin{enumerate}
    \item \textbf{Equal Weight ($1/N$):} A naive diversification baseline where capital is divided equally ($w_{i,t} = 0.20$ for all $N=5$ assets) and dynamically rebalanced daily.
    \item \textbf{Buy and Hold:} A static baseline where initial capital is equally divided on day 1 ($t=0$) and allowed to drift unhedged without subsequent rebalancing, incurring exactly zero ongoing transaction fees.
    \item \textbf{Markowitz Mean-Variance Optimization (MVO):} A dynamic quantitative benchmark that solves the constrained Sharpe ratio maximization problem on each trading day $t$ using an expanding historical window:
    \begin{align}
        \max_{\mathbf{w}_t} \quad & \frac{\mathbf{w}_t^\top \hat{\boldsymbol{\mu}}_t - R_f}{\sqrt{\mathbf{w}_t^\top \hat{\boldsymbol{\Sigma}}_t \mathbf{w}_t}} \\
        \text{s.t.} \quad & \sum_{i=1}^N w_{i,t} = 1.0, \quad 0 \le w_{i,t} \le 0.35, \; \forall i
    \end{align}
    solved numerically via Sequential Least Squares Programming (SLSQP).
    \item \textbf{PPO (No Regime Conditioned) [Ablation]:} An identical PPO network architecture trained with the regime indicators excluded ($D=16$), isolating the contribution of regime awareness.
\end{enumerate}

\subsection{Hyperparameter Configuration}
Both PPO models were implemented using Stable-Baselines3 and PyTorch. The policy and value networks employ a Multi-Layer Perceptron (MLP) architecture featuring two fully-connected hidden layers with 64 units each and $\tanh$ activation functions. Detailed hyperparameters are summarized in Table \ref{tab:hyperparams}.

\begin{table}[htbp]
\caption{PPO Agent Architecture and Training Hyperparameters}
\label{tab:hyperparams}
\centering
\begin{tabular}{ll}
\toprule
\textbf{Hyperparameter} & \textbf{Configured Value} \\
\midrule
Policy Architecture & MlpPolicy ($2 \times 64$ hidden units) \\
Activation Function & Tanh \\
Optimizer & Adam \\
Learning Rate ($\alpha$) & $3.0 \times 10^{-4}$ \\
Discount Factor ($\gamma$) & 0.99 \\
GAE Parameter ($\lambda$) & 0.95 \\
Clipping Parameter ($\epsilon$) & 0.20 \\
Entropy Coefficient ($c_2$) & 0.01 \\
Value Function Coefficient ($c_1$) & 0.50 \\
Batch Size & 64 \\
Rollout Horizon ($n_{\text{steps}}$) & 256 \\
Optimization Epochs ($n_{\text{epochs}}$) & 10 \\
Total Training Timesteps & 30,000 \\
Transaction Cost Rate ($c_{\text{fee}}$) & 10 bps ($0.0010$) \\
Single Asset Maximum Cap ($\omega_{\max}$) & 35\% ($0.35$) \\
Initial Virtual Portfolio Capital & INR 1,000,000 \\
\bottomrule
\end{tabular}
\end{table}

\subsection{Performance Evaluation Metrics}
Model evaluation is conducted across eight standard quantitative finance metrics:
\begin{enumerate}
    \item \textbf{Cumulative Return (CR):} Total percentage capital growth over the 502-day test period:
    \begin{equation}
        \text{CR} = \frac{V_T - V_0}{V_0} \times 100\%
    \end{equation}
    \item \textbf{Compound Annual Growth Rate (CAGR):}
    \begin{equation}
        \text{CAGR} = \left(\frac{V_T}{V_0}\right)^{\frac{252}{T}} - 1
    \end{equation}
    \item \textbf{Annualized Volatility ($\sigma_{\text{ann}}$):} Standard deviation of daily portfolio returns scaled to annual frequency:
    \begin{equation}
        \sigma_{\text{ann}} = \sqrt{252} \times \sqrt{\frac{1}{T-1} \sum_{t=1}^T (r_{\text{net}, t} - \bar{r})^2}
    \end{equation}
    \item \textbf{Sharpe Ratio (SR):} Risk-adjusted excess return per unit of total risk (assuming risk-free rate $R_f = 0$):
    \begin{equation}
        \text{SR} = \frac{\sqrt{252} \cdot \bar{r}_{\text{net}}}{\sigma_{\text{ann}}}
    \end{equation}
    \item \textbf{Maximum Drawdown (MDD):} Peak-to-trough decline in portfolio equity:
    \begin{equation}
        \text{MDD} = \max_{t \in [1, T]} \left( \frac{\max_{s \le t} V_s - V_t}{\max_{s \le t} V_s} \right) \times 100\%
    \end{equation}
    \item \textbf{Calmar Ratio:} Ratio of CAGR to Maximum Drawdown:
    \begin{equation}
        \text{Calmar} = \frac{\text{CAGR}}{\text{MDD}}
    \end{equation}
    \item \textbf{Total Turnover:} Cumulative portfolio reallocation volume:
    \begin{equation}
        \text{Turnover} = \sum_{t=1}^T \tau_t
    \end{equation}
    \item \textbf{Total Incurred Transaction Costs:} Cumulative friction in local currency (INR) given initial capital $V_0 = \text{INR } 1,000,000$:
    \begin{equation}
        \text{Cost} = \sum_{t=1}^T c_{\text{fee}} \cdot V_{t-1} \cdot \tau_t
    \end{equation}
\end{enumerate}

\subsection{Object-Oriented Virtual Execution Simulator}
To validate real-world deployability, we engineered a standalone, modular virtual execution engine in Java 17. The simulator implements an Object-Oriented architecture comprising `Asset`, `Order`, `Trade`, `Transaction`, `Portfolio`, `Account`, and `RiskManager` modules. Orders emitted by the Python policy bridge are transmitted to the Java execution server via REST sockets, where trade validation, lot sizing, cash ledger debits, and post-trade compliance checks are performed with microsecond-level latency.

\section{Empirical Results and Discussion}

\subsection{Benchmark Performance Comparison}
The out-of-sample backtest results across all five asset allocation strategies over the 502 unseen trading days (2024--2025) are presented in Table \ref{tab:results}.

\begin{table*}[t]
\caption{Out-of-Sample Performance Comparison on Unseen Test Dataset (2024--2025, 502 Trading Days)}
\label{tab:results}
\centering
\begin{tabular}{lcccccccc}
\toprule
\textbf{Strategy} & \textbf{Cumulative} & \textbf{CAGR} & \textbf{Sharpe} & \textbf{Annual} & \textbf{Max} & \textbf{Calmar} & \textbf{Total} & \textbf{Transaction Costs} \\
 & \textbf{Return (\%)} & \textbf{(\%)} & \textbf{Ratio} & \textbf{Volatility (\%)} & \textbf{Drawdown (\%)} & \textbf{Ratio} & \textbf{Turnover} & \textbf{(INR)} \\
\midrule
Equal Weight ($1/N$) & 102.20\% & 42.50\% & \textbf{1.46} & \textbf{25.00\%} & \textbf{26.48\%} & \textbf{1.60} & 6.11 & 8,696.21 \\
Buy \& Hold & \textbf{109.95\%} & \textbf{45.22\%} & 1.39 & 28.18\% & 28.69\% & 1.58 & \textbf{1.00} & \textbf{1,000.00} \\
Markowitz MVO & 106.17\% & 43.90\% & 1.33 & 29.01\% & 31.68\% & 1.39 & 21.45 & 30,984.32 \\
\midrule
\textbf{PPO (No Regime) [Ablation]} & 54.54\% & 24.47\% & 0.89 & 26.24\% & 32.22\% & 0.76 & 7.10 & 8,106.46 \\
\textbf{Regime-Aware PPO [Proposed]} & \textbf{99.80\%} & \textbf{41.64\%} & \textbf{1.33} & 27.52\% & 32.56\% & 1.28 & 53.26 & 70,390.37 \\
\bottomrule
\end{tabular}
\end{table*}

\subsection{Ablation Study: Critical Impact of Regime Conditioning}
The empirical results reveal a profound divergence between the regime-conditioned PPO agent and the unconditioned PPO baseline, validating our central research hypothesis:
\begin{itemize}
    \item \textbf{Return Enhancement:} Conditioning the PPO agent on market regimes increased cumulative return from \textbf{54.54\%} to \textbf{99.80\%}, representing an absolute gain of \textbf{+45.26\%} and an 83\% relative performance boost.
    \item \textbf{Risk-Adjusted Efficiency:} The Sharpe ratio improved substantially from \textbf{0.89} to \textbf{1.33} (+0.44), while the Calmar ratio surged from \textbf{0.76} to \textbf{1.28} (+68.4\%).
\end{itemize}

Without regime indicators, the agent's observation space lacks explicit structural signals differentiating secular bull phases from market turbulence. Consequently, the unconditioned policy converged to an overly conservative, sub-optimal allocation that failed to capture upside momentum during the strong 2024 tech rally, while still suffering comparable drawdown (32.22\%) during pullbacks. Conversely, the Regime-Aware PPO successfully learned state-contingent behavior: aggressively expanding allocations into high-beta winners (e.g., NVDA, MSFT) during `BULL` regimes, while rotating toward balanced defensive allocations when `HIGH\_VOLATILITY` states were triggered.

\subsection{Comparison with Classical Quantitative Benchmarks}
\begin{enumerate}
    \item \textbf{Equal Weight ($1/N$):} Equal Weight achieved a Sharpe ratio of 1.46 with a maximum drawdown of only 26.48\%, reaffirming DeMiguel et al.'s findings \cite{demiguel2009optimal} that simple equal allocation serves as a formidable out-of-sample benchmark due to absence of estimation error. However, naive equal weighting cannot dynamically adapt to regime shifts or capitalize on asset-specific momentum trends.
    \item \textbf{Buy \& Hold:} Buy and Hold generated a 109.95\% return with only INR 1,000 in transaction costs. However, this high return was driven almost entirely by unhedged concentration drift into NVDA as its market valuation skyrocketed during the 2024--2025 AI expansion. While profitable in a secular bull market, this passive drift exposes the portfolio to severe concentration risk, violating the 35\% regulatory cap enforced in our proposed active system.
    \item \textbf{Markowitz MVO:} Markowitz MVO achieved a cumulative return of 106.17\% and a Sharpe ratio of 1.33. However, MVO incurred significant rebalancing friction (INR 30,984.32 across 21.45 turnover units) due to rapid shifts in estimated sample covariance matrices. The Regime-Aware PPO achieved an identical Sharpe ratio of 1.33 while operating under direct continuous policy control without requiring numerical quadratic programming matrix inversions at inference time.
\end{enumerate}

\subsection{Transaction Cost and Turnover Analysis}
A common criticism of continuous reinforcement learning in finance is excessive churn. In our experiments, the Regime-Aware PPO incurred a total turnover of 53.26, resulting in INR 70,390.37 in cumulative transaction costs under the 10 bps fee model. Despite these realistic trading frictions, the model still achieved a net return of 99.80\%, nearly doubling initial investor capital over two years. 

Furthermore, the turnover penalty $\lambda_{\text{cost}} \tau_t$ in our multi-objective reward prevented the policy from destabilizing: without this turnover penalty, preliminary trial policies exhibited turnovers exceeding 180, destroying over 20\% of capital solely in transaction fees. This demonstrates that multi-objective turnover penalization is strictly necessary for real-world DRL deployment.

\section{Conclusion and Future Work}
In this paper, we introduced a robust, production-grade framework for multi-objective portfolio optimization combining continuous Proximal Policy Optimization (PPO) with deterministic market regime conditioning. By incorporating rolling asset momentum, trailing volatility, current portfolio inertia, and a 3-state macro regime classifier into a 19-dimensional MDP state space, the agent dynamically adjusts its asset allocations across shifting economic environments. 

Through rigorous out-of-sample testing on 502 unseen trading days (2024--2025) across five leading technology assets, the proposed Regime-Aware PPO achieved a cumulative return of 99.80\% (CAGR 41.64\%) and a Sharpe ratio of 1.33 net of transaction costs. Crucially, ablation experiments confirmed that regime conditioning provided an absolute return gain of +45.26\% and elevated the Sharpe ratio by +0.44 over an unconditioned DRL agent, confirming that regime awareness resolves policy paralysis during turbulent market phases.

Future extensions of this work include:
\begin{enumerate}
    \item Integrating Transformer-based temporal attention architectures (e.g., Temporal Fusion Transformers) to capture long-term sequential dependencies in place of feedforward MLPs.
    \item Incorporating macroeconomic NLP signals (e.g., Federal Reserve FOMC transcripts, financial news sentiment) into the regime detection engine.
    \item Scaling the asset universe to multi-asset classes, including fixed income, commodities, and sovereign bonds, to provide defensive diversification during structural bear regimes.
\end{enumerate}

\section*{Acknowledgment}
The authors express their deepest gratitude to their project supervisor, Ms. Vishakha Agarwal, Assistant Professor, Department of Computer Science and Engineering, School of Computer Science and Engineering (SCSE), IILM University, Greater Noida, for her invaluable mentorship, technical insights, and continuous guidance throughout the duration of this research. We also thank the SCSE faculty for providing the computing resources and academic environment to carry out this work.

\bibliographystyle{IEEEtran}
\begin{thebibliography}{00}

\bibitem{markowitz1952portfolio}
H.~Markowitz, ``Portfolio selection,'' \emph{The Journal of Finance}, vol.~7, no.~1, pp. 77--91, 1952.

\bibitem{mandelbrot1963variation}
B.~Mandelbrot, ``The variation of certain speculative prices,'' \emph{The Journal of Business}, vol.~36, no.~4, pp. 394--419, 1963.

\bibitem{cont2001empirical}
R.~Cont, ``Empirical properties of asset returns: Stylized facts and statistical issues,'' \emph{Quantitative Finance}, vol.~1, no.~2, pp. 223--236, 2001.

\bibitem{michaud1989markowitz}
R.~O. Michaud, ``The Markowitz optimization enigma: Is 'optimized' optimal?'' \emph{Financial Analysts Journal}, vol.~45, no.~1, pp. 31--42, 1989.

\bibitem{sharpe1964capital}
W.~F. Sharpe, ``Capital asset prices: A theory of market equilibrium under conditions of risk,'' \emph{The Journal of Finance}, vol.~19, no.~3, pp. 425--442, 1964.

\bibitem{lintner1965valuation}
J.~Lintner, ``The valuation of risk assets and the selection of risky investments in stock portfolios and capital budgets,'' \emph{The Review of Economics and Statistics}, vol.~47, no.~1, pp. 13--37, 1965.

\bibitem{black1992global}
F.~Black and R.~Litterman, ``Global portfolio optimization,'' \emph{Financial Analysts Journal}, vol.~48, no.~5, pp. 28--43, 1992.

\bibitem{ledoit2004well}
O.~Ledoit and M.~Wolf, ``A well-conditioned estimator for large-dimensional covariance matrices,'' \emph{Journal of Multivariate Analysis}, vol.~88, no.~2, pp. 365--411, 2004.

\bibitem{demiguel2009optimal}
V.~DeMiguel, L.~Garlappi, and R.~Uppal, ``Optimal versus naive diversification: How inefficient is the 1/N portfolio strategy?'' \emph{The Review of Financial Studies}, vol.~22, no.~5, pp. 1915--1953, 2009.

\bibitem{moody2001learning}
J.~Moody and M.~Saffell, ``Learning to trade via direct reinforcement,'' \emph{IEEE Transactions on Neural Networks}, vol.~12, no.~4, pp. 875--889, 2001.

\bibitem{jiang2017deep}
Z.~Jiang, D.~Xu, and J.~Liang, ``A deep reinforcement learning framework for the financial portfolio management problem,'' \emph{arXiv preprint arXiv:1706.10059}, 2017.

\bibitem{liang2018adversarial}
Z.~Liang, H.~Chen, J.~Zhu, K.~Jiang, and Y.~Li, ``Adversarial deep reinforcement learning in portfolio management,'' \emph{arXiv preprint arXiv:1808.09940}, 2018.

\bibitem{ding2020deep}
K.~Ding, Y.~Lv, and C.~Zheng, ``Deep reinforcement learning for portfolio management: A multi-agent approach,'' in \emph{Proceedings of the IEEE International Conference on Data Mining (ICDM)}, 2020, pp. 112--120.

\bibitem{zhang2020deep}
Z.~Zhang, S.~Zohren, and S.~Roberts, ``Deep reinforcement learning for trading,'' \emph{The Journal of Financial Data Science}, vol.~2, no.~2, pp. 25--40, 2020.

\bibitem{hamilton1989new}
J.~D. Hamilton, ``A new approach to the economic analysis of nonstationary time series and the business cycle,'' \emph{Econometrica}, vol.~57, no.~2, pp. 357--384, 1989.

\bibitem{guidolin2011markov}
M.~Guidolin, ``Markov switching models in empirical finance,'' in \emph{Missing Data Methods: Time-Series Methods and Applications}, Emerald Group Publishing Limited, 2011, pp. 1--86.

\bibitem{ang2002regime}
A.~Ang and G.~Bekaert, ``Regime switches in interest rates,'' \emph{Journal of Business \& Economic Statistics}, vol.~20, no.~2, pp. 163--182, 2002.

\bibitem{kritzman2012regime}
M.~Kritzman, S.~Page, and D.~Turkington, ``Regime shifts: Implications for asset allocation,'' \emph{Financial Analysts Journal}, vol.~68, no.~3, pp. 22--39, 2012.

\bibitem{schulman2017proximal}
J.~Schulman, F.~Wolski, P.~Dhariwal, A.~Radford, and O.~Klimov, ``Proximal policy optimization algorithms,'' \emph{arXiv preprint arXiv:1707.06347}, 2017.

\bibitem{schulman2015high}
J.~Schulman, P.~Moritz, S.~Levine, M.~Jordan, and P.~Abbeel, ``High-dimensional continuous control using generalized advantage estimation,'' \emph{arXiv preprint arXiv:1506.02438}, 2015.

\end{thebibliography}

\end{document}
